
\documentclass{article}
\usepackage{colortbl}
\usepackage{makecell}
\usepackage{multirow}
\usepackage{supertabular}

\begin{document}

\newcounter{utterance}

\centering \large Interaction Transcript for game `cryptolect', experiment `basic', episode 6 with BSU{-}SLIM\_\_Qwen3.5{-}9B{-}Base\_\_playpen{-}prm{-}9b{-}best\_of\_n.
\vspace{24pt}

{ \footnotesize  \setcounter{utterance}{1}
\setlength{\tabcolsep}{0pt}
\begin{supertabular}{c@{$\;$}|p{.15\linewidth}@{}p{.15\linewidth}p{.15\linewidth}p{.15\linewidth}p{.15\linewidth}p{.15\linewidth}}
    \# & \multicolumn{2}{c}{Player} && \multicolumn{2}{c}{Game Master} \\
    \hline

    \theutterance \stepcounter{utterance}  
    & & & \multicolumn{4}{p{0.6\linewidth}}{
        \cellcolor[rgb]{0.9,0.9,0.9}{
            \makecell[{{p{\linewidth}}}]{
                \texttt{\tiny{[P1$\langle$GM]}}
                \texttt{You are a field linguist. You have made contact with the only speaker of an unknown} \\
\texttt{language (the "informant"). Nobody has ever written down this language before. Your task} \\
\texttt{is to work out how it works and then translate sentences into it correctly.} \\
\\ 
\texttt{You will study the language by asking the informant to translate English sentences for} \\
\texttt{you. You may only talk about a small fixed world, using exactly these English words:} \\
\\ 
\texttt{Nouns: dog, bird, stone, river, hunter, king, road, ship} \\
\texttt{Adjectives: (this language has no adjectives)} \\
\texttt{Verbs (write them in this base form): feed, see, hear, hold, watch} \\
\texttt{Articles: the (definite), a (indefinite)} \\
\texttt{Number words: two, three (no number word = one / singular)} \\
\texttt{These are ALL the words in play {-} the language has no other words.} \\
\\ 
\texttt{HOW TO WRITE AN ENGLISH SENTENCE} \\
\texttt{Every sentence is one clause with a subject, a verb, and an object, written in this fixed English notation:} \\
\texttt{    SUBJECT  verb  OBJECT} \\
\texttt{Each of SUBJECT and OBJECT is a noun phrase of the form:} \\
\texttt{    {[}the | a{]}  {[}number word{]}  {[}adjective ...{]}  noun} \\
\texttt{{-} 'the' = definite, 'a' = indefinite (singular only); leave the article off for an indefinite plural.} \\
\texttt{{-} a number word (two, three) sets the count; with no number word the noun is singular.} \\
\texttt{Examples of well{-}formed English sentences you may send:} \\
\texttt{    a dog feed a bird} \\
\texttt{    three dogs feed a bird} \\
\\ 
\texttt{WHAT YOU ALREADY KNOW} \\
\texttt{The informant has already translated these sentences for you. Each line is} \\
\texttt{"<English>  =>  <the language>":} \\
\texttt{    three birds feed a bird  =>  vevogo vi vevo pove} \\
\texttt{    a river hear three hunters  =>  vopotu rugigo vi guvi} \\
\texttt{    three stones hold three kings  =>  kevogo vi bopigo vi rite} \\
\texttt{    two roads hear two dogs  =>  gavogago bete gepego bete guvi} \\
\texttt{    a bird feed a dog  =>  vevo gepe pove} \\
\texttt{    a river see three birds  =>  vopotu vevogo vi tugige} \\
\texttt{    three stones watch a ship  =>  kevogo vi tipo guku} \\
\\ 
\texttt{HOW THE SESSION WORKS} \\
\texttt{1. INVESTIGATION: You have 4 rounds of questioning. In each round you may ask} \\
\texttt{   the informant to translate up to 3 English sentences (so at most} \\
\texttt{   12 probes in total). The informant answers each one truthfully and} \\
\texttt{   exactly. Use these answers to figure out the words and the grammar (word order, how} \\
\texttt{   number/case/definiteness are marked, how words must agree, and any sound changes).} \\
\texttt{   Choose your probes deliberately {-} test specific hypotheses.} \\
\texttt{2. EXAM: After the investigation rounds, the informant will give you 10 NEW} \\
\texttt{   English sentences. You must translate each one INTO the language. Every answer is} \\
\texttt{   graded by EXACT match against how the informant would say it {-} every word, ending and} \\
\texttt{   space must be right (capitalisation and surrounding spaces are ignored). The exam} \\
\texttt{   sentences will use the same vocabulary but in combinations you have not been shown, so} \\
\texttt{   you must induce the general rules, not just memorise examples.} \\
\\ 
\texttt{RESPONSE FORMAT (for every investigation round)} \\
\texttt{List your probe sentences under a 'PROBES:' line, one per line. A THOUGHTS line is} \\
\texttt{optional {-} keep any reasoning brief so you always finish writing your PROBES:} \\
\texttt{THOUGHTS: <optional; brief>} \\
\texttt{PROBES:} \\
\texttt{1. <an English sentence, written as described above>} \\
\texttt{2. <another English sentence>} \\
\texttt{(up to 3 lines; at least one)} \\
\\ 
\texttt{Begin your first round of probes now.} \\
            }
        }
    }
     \\ \\

    \theutterance \stepcounter{utterance}  
    & \multicolumn{4}{p{0.6\linewidth}}{
        \cellcolor[rgb]{0.561,0.737,0.561}{
            \makecell[{{p{\linewidth}}}]{
                \texttt{\tiny{[P1$\rangle$GM]}}
                \texttt{PROBES:} \\
\texttt{1. a dog hear a bird} \\
\texttt{2. a dog hear two birds} \\
\texttt{3. a dog hear three birds} \\
\\ 
\texttt{THOUGHTS:} \\
\texttt{The initial data shows a clear word order of SVO (Subject{-}Verb{-}Object).} \\
\texttt{I need to determine how the languages marks specific grammatical features.} \\
\texttt{Comparing `three birds feed a bird` (vevogo vi vevo pove) with `a bird feed a dog` (vevo gepe pove):} \\
\texttt{{-} The verb `feed` maps to `pove` in both cases, suggesting the verb form is invariant (or depends only on tense, which isn't here).} \\
\texttt{{-} The object changes: `a bird` {-}> `pove`, `a dog` {-}> `gepe`. So `bird` = `pove`, `dog` = `gepe`.} \\
\texttt{{-} The subject changes: `a bird` {-}> `vevogo`, `a dog` {-}> `vevogo`. Wait, `a bird` is `vevo` in the second example, but `vevogo` in the first? This suggests that the subject marking might depend on the number or the specific noun, OR I misidentified the mapping immediately.} \\
\texttt{Let's re{-}examine:} \\
\texttt{`three birds feed a bird` => `vevogo vi vevo pove`} \\
\texttt{Subject: `three birds` {-}> `vevogo`} \\
\texttt{Object: `a bird` {-}> `pove`} \\
\texttt{`a bird feed a dog` => `vevo gepe pove`} \\
\texttt{Subject: `a bird` {-}> `vevo`} \\
\texttt{Object: `a dog` {-}> `gepe`} \\
\texttt{This confirms `bird` = `pove` and `dog` = `gepe`.} \\
\texttt{The difference between `vevogo` (three birds) and `vevo` (a bird) suggests that number is marked on the subject.} \\
\texttt{Also, looking at `a river see three birds` => `vopotu vevogo vi tugige`:} \\
\texttt{Subject: `a river` {-}> `vopotu` (Wait, `three rivers` wasn't in the list, but `a river` is here).} \\
\texttt{Object: `three birds` {-}> `vevogo vi tugige`? No, the object phrase is `vevogo vi tugige`.} \\
\texttt{In `three birds feed a bird`, the object was just `pove`.} \\
\texttt{Here the object is `three birds`, which seems to map to `vevogo vi tugige`.} \\
\texttt{This implies `three birds` might be a compound or a specific form involving the singular form plus a marker.} \\
\texttt{Hypothesis: `vi` acts as a plural marker or a specific number marker for "three"?} \\
\texttt{Let's test number marking on the subject and object separately.} \\
\texttt{Subject: `three birds` = `vevogo`. Subject: `a bird` = `vevo`.} \\
\texttt{Hypothesis: Number is marked by adding a suffix `{-}o` for three? `vevo` + `go`? Or `vi`?} \\
\texttt{Actually, `vevogo` looks like `vevo` + `go`.} \\
\texttt{But in the object of `three birds` in the sentence `a river see three birds`, the object form is `vevogo vi tugige`.} \\
\texttt{If `bird` is `pove`, then `three birds` (object) should be `pove` plus some marker.} \\
\texttt{But it is `vevogo vi tugige`. This is very strange if `bird` is `pove`.} \\
\texttt{Wait, let me re{-}read the initial data carefully.} \\
\texttt{`three birds feed a bird` => `vevogo vi vevo pove`} \\
\texttt{`a river see three birds` => `vopotu vevogo vi tugige`} \\
\texttt{In the second sentence, the object is `three birds`. The translation is `vevogo vi tugige`.} \\
\texttt{In the first sentence, the object is `a bird`. The translation is `pove`.} \\
\texttt{So `three birds` != `pove`. It is `vevogo vi tugige`.} \\
\texttt{This suggests that `vevogo` and `tugige` are both part of the phrase for "three birds".} \\
\texttt{Maybe `vevogo` is the marker for "bird" in the plural/three context, and `tugige` is another marker?} \\
\texttt{Or maybe `bird` isn't `pove`?} \\
\texttt{Let's check `a river feed three birds`? I don't have that directly.} \\
\texttt{Let's check `three stones watch a ship` => `kevogo vi tipo guku`.} \\
\texttt{Object: `a ship` => `guku`.} \\
\texttt{Subject: `three stones` => `kevogo vi`.} \\
\texttt{Compare `a river` (`vopotu`) and `three stones` (`kevogo vi`).} \\
\texttt{Compare `a bird` (`vevo`) and `three birds`} \\
            }
        }
    }
    & & \\ \\

    \theutterance \stepcounter{utterance}  
    & & & \multicolumn{2}{p{0.3\linewidth}}{
        \cellcolor[rgb]{0.9,0.9,0.9}{
            \makecell[{{p{\linewidth}}}]{
                \texttt{\tiny{[GM$|$GM]}}
                \texttt{at most 3 probes per round, but 48 were given} \\
            }
        }
    }
    & & \\ \\

    \theutterance \stepcounter{utterance}  
    & & & \multicolumn{4}{p{0.6\linewidth}}{
        \cellcolor[rgb]{0.9,0.9,0.9}{
            \makecell[{{p{\linewidth}}}]{
                \texttt{\tiny{[P1$\langle$GM]}}
                \texttt{Your last response could not be processed: at most 3 probes per round, but 48 were given.} \\
\texttt{List your probes (a brief THOUGHTS line is optional):} \\
\texttt{PROBES:} \\
\texttt{1. <an English sentence>} \\
\texttt{(up to 3 probe lines; at least one)} \\
            }
        }
    }
     \\ \\

    \theutterance \stepcounter{utterance}  
    & \multicolumn{4}{p{0.6\linewidth}}{
        \cellcolor[rgb]{0.561,0.737,0.561}{
            \makecell[{{p{\linewidth}}}]{
                \texttt{\tiny{[P1$\rangle$GM]}}
                \texttt{PROBES:} \\
\texttt{1. a dog hear a bird} \\
\texttt{2. a dog hear two birds} \\
\texttt{3. a dog hear three birds} \\
\\ 
\texttt{THOUGHTS:} \\
\texttt{I need to establish the base forms for the nouns (dog, bird) and the specific markers for different numbers (two, three) and definiteness. The initial data shows discrepancies between singular and plural markers that depend on the number (e.g., "a bird" vs "three birds"). I will start by probing the interaction of the verb "hear" with "dog" and "bird" across different numbers and definiteness to isolate the affixes. This includes checking if the subject or the object carries the number/definiteness markers, and how "two" differs from "three".} \\
            }
        }
    }
    & & \\ \\

    \theutterance \stepcounter{utterance}  
    & & & \multicolumn{2}{p{0.3\linewidth}}{
        \cellcolor[rgb]{0.9,0.9,0.9}{
            \makecell[{{p{\linewidth}}}]{
                \texttt{\tiny{[GM$|$GM]}}
                \texttt{at most 3 probes per round, but 5 were given} \\
            }
        }
    }
    & & \\ \\

\end{supertabular}
}

\end{document}
